%% =====================================================================
%%  Corrigés — Unit 17 « Where were you? »
%% =====================================================================
\corrigesUnite{17}{Where were you?}

\begin{corrige}{1}
\begin{reponses}[4]
\item \rep{was}
\item \rep{were}
\item \rep{was}
\item \rep{were}
\item \rep{was}
\item \rep{were}
\item \rep{was}
\item \rep{were}
\end{reponses}
\pourquoi{un nom au singulier (\en{Emma}, \en{the party}) se comporte comme \en{she} ou \en{it} :
\en{was}. Un pluriel (\en{Léa's parents}, \en{Jack and Kate}) comme \en{they} : \en{were}.}
\piege{\en{you} prend toujours \en{were}, même pour une seule personne : \en{You were late}.}
\end{corrige}

\begin{corrige}{2}
\begin{reponses}[2]
\item \rep{She was at a concert.}
\item \rep{He was in bed.}
\item \rep{She was at the airport.}
\item \rep{He was at the café.}
\item \rep{He was at the office.}
\item \rep{She was at the supermarket.}
\end{reponses}
\retiens{\en{at home}, \en{at work}, \en{at the office}, mais \en{in bed} : pas d'article dans ces
expressions toutes faites.}
\end{corrige}

\begin{corrige}{3}
\begin{reponses}[2]
\item \rep{Yes, I was.}
\item \rep{No, he wasn't.}
\item \rep{Yes, they were.}
\item \rep{No, it wasn't.}
\item \rep{Yes, she was.}
\item \rep{No, we weren't.}
\end{reponses}
\pourquoi{la réponse courte reprend le sujet sous forme de pronom (\en{Jack} → \en{he}, \en{the weather}
→ \en{it}) et \en{was} ou \en{were}. Question 6 : \en{you and Sam} → \en{we}.}
\end{corrige}

\begin{corrige}{4}
\begin{reponses}[3]
\item \rep{there was}
\item \rep{there were}
\item \rep{There wasn't}
\item \rep{there weren't}
\item \rep{there were}
\item \rep{there was}
\end{reponses}
\pourquoi{\en{there was} + singulier (\en{a cinema}, \en{music}) ; \en{there were} + pluriel
(\en{two bakeries}, \en{children}). Avec \en{any}, on emploie la négation : \en{there weren't any}.}
\end{corrige}

\begin{corrige}{5}
\begin{reponses}[3][french]
\item \rep{V}
\item \rep{V}
\item \rep{F}
\item \rep{F}
\item \rep{V}
\item \rep{F}
\end{reponses}
\pourquoi{le pont de Brooklyn date de 1883 ; la statue de la Liberté est un cadeau de la France ;
l'Empire State Building n'existe que depuis 1931.}
\end{corrige}

\begin{corrige}{6}
\begin{reponses}[1]
\item \rep{I was at the party.}
\item \rep{They were very tired.}
\item \rep{Were you at work yesterday?}
\item \rep{I was born in Lyon.}
\item \rep{She was here two days ago.}
\end{reponses}
\piege{phrase 4 : « je suis né » est au passé en anglais : \en{I was born}. Phrase 5 : « il y a deux
jours » = \en{two days ago}, jamais \en{there is two days}.}
\end{corrige}

\begin{corrige}{7}
\begin{reponses}[3]
\item \rep{was}
\item \rep{were}
\item \rep{were}
\item \rep{was}
\item \rep{was}
\item \rep{were}
\end{reponses}
\pourquoi{dans une phrase, \en{was} se réduit à \phon{wəz} (on entend un petit \phon{z}) et \en{were} à
\phon{wɚ} (on entend un \en{r}). Aidez-vous aussi du sujet ! Texte de la piste
page~\pageref{trans-17.10}.}
\end{corrige}

\begin{corrige}{8}
\begin{reponses}[2]
\item \rep{kitchen}
\item \rep{post office}
\item \rep{client}
\item \rep{desk}
\item \rep{wasn't}
\end{reponses}
\pourquoi{Léa dit que Jack était avec elle, à son bureau : Jack n'était donc pas à la poste… et le
gâteau était délicieux ! Texte de la piste page~\pageref{trans-17.11}.}
\end{corrige}

\begin{corrige}{9}
\begin{reponses}[4]
\item \rep{was}
\item \rep{were}
\item \rep{Were}
\item \rep{were}
\item \rep{cold}
\item \rep{where}
\item \rep{were}
\item \rep{ago}
\end{reponses}
\end{corrige}

\begin{corrige}{10}
\begin{reponses}[1]
\item \rep{Where were you last night? — I was at home.}
\item \rep{How was the concert? — It was great!}
\item \rep{They weren't late.}
\item \rep{Were you tired? — Yes, I was.}
\item \rep{There were a lot of people.}
\item \rep{I was born in Lyon thirty years ago.}
\item \rep{There wasn't a supermarket.}
\item \rep{Last week, it was cold.}
\end{reponses}
\piege{phrase 5 : « il y avait beaucoup de monde » → \en{There were a lot of people} : \en{people} est
un pluriel.}
\end{corrige}

\begin{corrige}{11}
Réponse modèle :
\begin{anglais}
\begin{tcolorbox}[otBase,colback=otGreenTint,colframe=otGreenTint,fontupper=\fontsize{10.4}{14}\selectfont]
I was born in Nantes in 1995. Our house was small, but there was a big garden. There was a park and
a swimming pool near our street. My school was next to the church. I was very happy there!
\end{tcolorbox}
\end{anglais}
\retiens{\en{a swimming pool} = « une piscine » ; \en{a church} = « une église ».}
\end{corrige}

\begin{corrige}{12}
Réponses modèles :
\begin{reponses}[1]
\item \rep{I was at home.} \textit{ou} \rep{I was at a friend's house.}
\item \rep{It was great!} \textit{ou} \rep{It was very quiet.}
\item \rep{Yes, I was.} \textit{ou} \rep{No, I wasn't.}
\item \rep{I was born in Toulouse, in 1990.}
\item \rep{Yes, there was. It was next to the station.}
\end{reponses}
\end{corrige}

\begin{corrigetest}
\begin{reponses}[5][french]
\item \rep{a}
\item \rep{b}
\item \rep{c}
\item \rep{a}
\item \rep{b}
\item \rep{c}
\item \rep{a}
\item \rep{c}
\item \rep{b}
\item \rep{a}
\end{reponses}
\pourquoi{question 4 : \en{There were} devant un pluriel ; \en{It had} est un calque du français.
Question 3 : la réponse courte reprend \en{were}, jamais \en{did}.}
\end{corrigetest}
